\section{The named residuals and outlook}
\label{sec:residual}

The program reduces Frankl, along each of its two arcs, to a single
precisely named open object. This section collects those objects, says
what closing each would buy, and records the recommended next steps. It
introduces no new mathematics.

\subsection{The two residuals}
\label{ssec:two-residuals}

\begin{description}[leftmargin=1.6em,style=nextline]

\item[The cohomological residual --- the $Y_1$ computation.]
By Theorem~\ref{thm:y1-reduction} and Open residual~\ref{res:y1},
closing the cohomological arc reduces to one homotopy-colimit
computation: the $\walsh S$-isotype cohomology
$\isotype Sk(\Xn/X(\Dx))$ of the cofiber of the matched-cylinder
inclusion, for $S\subsetneq[n]$ and $k\ge1$. Concretely this is the
cohomology of the bar strings of $\Cn$ that exit the matched
subcategory $\Dx$ --- equivalently, the comma-category twist
$N(\rho_x/-)$ of the trace functor $\rho_x$, the homotopy type of the
stuck-family fibres. If this cohomology vanishes, $Y_1$ holds,
Conjecture~\ref{conj:uc10-1} follows (with the top-Walsh half from
UC14~$R3$, whose method is sound), Fact Two
(Proposition~\ref{prop:fact-two}) is established, and the closing
contradiction of Section~\ref{ssec:closing} goes through. If it does
not vanish, the two-part contradiction in its current form fails and
the cohomological arc needs a new idea.

\item[The combinatorial residual --- the boundary inequality.]
By Theorem~\ref{thm:boundary-conservation} and
Section~\ref{ssec:combinatorial-status}, closing the combinatorial arc
reduces to one explicit integer inequality: the global sign of the
boundary deficit $\sum_R B(R)$ against the rooted mass
$\sum_R\sum_{S_0}a_{S_0}|S_0|$, together with the excess-co-fibre term
$\mathcal S(R\cup R')\setminus(\mathcal S(R)\vee\mathcal S(R'))$. Every
quantity is a finite, explicitly computable count;
Proposition~\ref{prop:no-per-fibre} shows it is irreducibly a
sum-level statement. Closing it proves Frankl outright; it is currently
settled only on the structured hierarchy of
Section~\ref{ssec:combinatorial-status}.

\end{description}

\begin{remark}[The two residuals are the same difficulty in two languages]
\label{rem:same-difficulty}
It is worth recording the program's own reading. The combinatorial
residual is a statement about the \emph{co-fibres} $\mathcal S(R)$ ---
union-closed families on the generator $W$ --- and their failure to
assemble monotonically across the fibre lattice $\mathcal R_W$. The
cohomological residual is a statement about the \emph{stuck-family
fibres} of the trace functor and their failure to be contractible. Both
are, in the end, about the same phenomenon flagged in
Remark~\ref{rem:asymmetry} and Proposition~\ref{prop:transport-identity}:
the obstruction to gluing the local coordinate symmetries into a global
one. Whether the two residuals are formally inter-reducible is itself
open and would be a natural object of study; the program has not
established a bridge between them.
\end{remark}

\subsection{What would, and would not, constitute a proof}
\label{ssec:what-closes}

For the avoidance of doubt:

\begin{itemize}[leftmargin=1.6em]
\item Closing the combinatorial residual --- proving the boundary
  inequality unconditionally --- would prove Frankl, with no
  dependence on the cohomological arc.
\item Closing the cohomological residual --- computing
  $\isotype Sk(\Xn/X(\Dx))$ and finding it zero --- would, \emph{together
  with} a referee-grade verification of the remaining inputs
  (Conjecture~\ref{conj:uc10-1}'s top-Walsh half via $R3$; the
  no-cup-product faithfulness check for $Y_2$,
  Remark~\ref{rem:source-level1}; and the UC11 Lemma~6.1 equivalence,
  Remark~\ref{rem:wrapper-risk}), prove Frankl.
\item Neither residual is closed. Re-banding an internal verdict from
  amber to green does not close anything; only the mathematics does.
\item The external input of Remark~\ref{rem:external-input} must be
  independently verified against its source before the cohomological
  arc can claim to be self-contained.
\end{itemize}

\subsection{Recommended next steps}
\label{ssec:next-steps}

Consistent with the audit (Section~\ref{sec:proof-state}) the program
record's recommended priorities are:

\begin{enumerate}[leftmargin=1.8em]
\item \textbf{Attack the $Y_1$ residual directly}
  (Open residual~\ref{res:y1}). The recommended entry point is to
  compute the comma categories $\rho_x/(T,\G)$ --- the fibres of the
  trace to $[n]\setminus\{x\}$ --- and their nerves; the stuck families
  are exactly the non-contractible-fibre locus. The smallest genuinely
  open instance is $n=3$, $S=\{1\}$, $k=1$, which should be computed by
  hand first. Do \emph{not} attempt a cofinality argument for $\rho_x$:
  by Correction~\ref{corr:r2} that is equivalent to the false $R2$.
\item \textbf{Do not re-attempt $Y_1$ by the UC13~\S4.5 route.} The
  trace-exactness error (Correction~\ref{corr:trace-exact}) is
  structural; the next genuine ticket is the residual, not another
  bridge-iteration argument.
\item \textbf{Do not cite UC14~$R2$, and do not cite UC14 or UC13 as
  ``green''.} $R2$ is false (Correction~\ref{corr:r2}); UC13's closing
  banner is withdrawn (Status~\ref{status:closing}).
\item \textbf{Treat the formalisation honestly.} The headline
  \texttt{Frankl\_Holds} is axiom-dependent and the axiom is
  Frankl-in-disguise (Proposition~\ref{prop:axiom-is-frankl}); the
  isotype encodings are answer-steered
  (Section~\ref{ssec:answer-steered}). A faithful formalisation of
  $Y_1$ should wait until $Y_1$ has a valid proof --- there is
  currently nothing valid to formalise.
\item \textbf{The combinatorial arc is independently publishable as
  partial progress.} Sections~\ref{sec:setup}--\ref{sec:combinatorial}
  --- the transport identity, the trace-partition theorem, the
  large-fibre localisation, the boundary conservation law, the
  no-per-fibre-sign impossibility result, and the settled structured
  hierarchy --- are a self-contained body of correct results that do
  not depend on any disputed step. They are the program's solid
  deliverable.
\end{enumerate}

\subsection{Summary}
\label{ssec:final-summary}

The program is a genuine, structured attack on a hard conjecture, and
it has produced real mathematics: a clean reduction of Frankl to an
explicit finitary inequality with a growing settled hierarchy; a native
recovery of the classical small-generator cases; a genuine
impossibility result; an equivariant cohomology theory of
intersection-closed families; one unconditional trace-injectivity
theorem; and a mechanism-sound isotype-landing result. It has not
proved Frankl, and the audit consolidated here corrects an earlier
internal over-claim to that effect: the closing contradiction depends
on the lower-Walsh vanishing $Y_1$, whose circulated proof is invalid
and whose advertised repair is false. What the program honestly stands
on today is two correct reductions to two precisely named open objects.
The recommended next move is the $Y_1$ residual of
Open residual~\ref{res:y1}.
